\documentclass{article}
\usepackage[utf8]{inputenc}
\usepackage{amsmath}
\usepackage{mathtools}
\usepackage[sexy]{evanchen} 
\usepackage{graphicx}
\graphicspath{{8 Outsource K Math/Gambar/}}

\title{LTE}
\author{Azzam and Akeyla}

\begin{document}

\maketitle
\begin{center}
    Lemma LTE (Lifting The Exponent Lemma) merupakan salah satu topik yang baru di teori bilangan dan merupakan metode yang cukup ampuh dalam menyelesaikan beberapa persamaan Diophantine eksponensial . \\
    Lemma LTE membantu kita dalam menemukan pangkat terbesar dari suatu bilangan prima $p$ yang membagi $a^n\pm b^n$ untuk suatu bilangan asli $a,b.$
\end{center}

\section{Definisi}
Definisikan $v_p(a)$ sebagai bilangan bulat terbesar sehingga $p^{v_p(a)}|a$ , dimana $p$ merupakan bilangan prima dan $a$ merupakan bilangan asli .
\\ Contoh : $v_3(54)=3$ , $v_5(100)=2$ , $v_2(2)=1$ , $v_5(7)=0$ .

\section{Dua Lemma yang cukup penting}

\begin{lemma}
Misalkan $x$ dan $y$ adalah bilangan bulat dan $n$ adalah bilangan asli . Diberikan sebarang bilangan prima $p$ yang memenuhi FPB$(n,p)=1$ , $p|x-y$ dan $x,y$ keduanya tidak habis dibagi oleh $p$ . Maka $$v_p(x^n-y^n)=v_p(x-y).$$
\end{lemma}
\begin{proof}[Bukti]
Amati bahwa $$x^n-y^n=(x-y)(x^{n-1}+x^{n-2}y+\dots +y^{n-1})$$
Maka , kita hanya perlu membuktikan kalau $(x^{n-1}+x^{n-2}y+\dots +y^{n-1})$ tidak habis dibagi $p$
Karena $x\equiv y$(mod $p$) , maka $$(x^{n-1}+x^{n-2}y+\dots +y^{n-1}) \equiv nx^n-1(mod p)$$
Akibatnya , $p$ tidak membagi $(x^{n-1}+x^{n-2}y+\dots +y^{n-1})$
dan $v_p(x^n-y^n)=v_p(x-y)$
\end{proof}

\begin{lemma}
Misalkan $x$ dan $y$ adalah bilangan bulat dan $n$ adalah bilangan ganjil . Diberikan sebarang bilangan prima $p$ yang memenuhi FPB$(n,p)=1$ , $p|x+y$ dan $x,y$ keduanya tidak habis dibagi oleh $p$ . Maka $$v_p(x^n+y^n)=v_p(x+y).$$
\end{lemma}

\begin{proof}[Bukti]
Dengan cara yang sama dengan Lemma 1 , kecuali bahwa $x \equiv -y(mod p)$
\end{proof} 

\section{Lemma LTE}

\begin{lemma}[LTE Tipe 1]
Misalkan $x$ dan $y$ adalah bilangan bulat dan $n$ adalah bilangan asli . Diberikan sebarang bilangan prima \textit{ganjil} $p$ yang memenuhi $p|x-y$ dan $x,y$ keduanya tidak habis dibagi oleh $p$ . Maka $$v_p(x^n-y^n)=v_p(x-y)+v_p(n).$$
\end{lemma}

\begin{lemma}[LTE Tipe 2]
Misalkan $x$ dan $y$ adalah bilangan bulat dan $n$ adalah bilangan ganjil positif. Diberikan sebarang bilangan prima \textit{ganjil} $p$ yang memenuhi $p|x+y$ dan $x,y$ keduanya tidak habis dibagi oleh $p$ . Maka $$v_p(x^n+y^n)=v_p(x+y)+v_p(n).$$
\end{lemma}

\begin{lemma}[LTE Tipe 3]
Misalkan $x$ dan $y$ adalah bilangan ganjil sehingga $4|x-y$ dan $n$ adalah bilangan asli . Maka $$v_2(x^n-y^n)=v_2(x-y)+v_2(n)$$
\end{lemma}

\begin{lemma}[LTE Tipe 4]
Misalkan $x$ dan $y$ adalah bilangan ganjil dan $n$ adalah bilangan ganjil positif . Maka $$v_2(x^n-y^n)=v_2(x-y)+v_2(x+y)+v_2(n)-1$$
\end{lemma}

\section{Contoh Penggunaan}
\begin{example}
Misal $p$ adalah bilangan prima yang lebih besar dari 2016 , $a,b$ adalah bilangan bulat sehingga $p|a+b , p^2|a+b$ dan $p^2|a^{2015}+b^{2015}$ . Ada berapa banyak nilai $n$ bilangan asli $\leq 2016$ sehingga $p^n|a^{2016}+b^{2016}$ ?
\end{example}
\begin{proof}[Solusi]
Asumsikan $a$ dan $b$ tidak habis dibagi oleh $p$.\\
Karena $p>2016$ , maka $p$ tidak membagi 2015 dan $p$ ganjil .\\
Dengan Lemma 2 , didapat $v_p(a^{2015}+b^{2015})=v_p(a+b)=1.$\\
Padahal , $p^2|a^{2015}+b^{2015}$ , kontradiksi .\\
Akibatnya $p|a$ dan $p|b$ . \\
Misalkan $a=xp$ dan $b=yp$ untuk suatu bilagan bulat $x,y$ .\\
Maka , $a^{2016}+b^{2016}=p^{2016}(x^{2016}+y^{2016})$.\\
Ada 2016 nilai $n$ yang memenuhi , yakni ${{1,2,3\dots 2016}}.$
\end{proof}
\begin{example}
Misalkan $a,n$ bilangan asli , $p$ bilangan prima ganjil dan memenuhi $a^p \equiv 1(mod p^n)$ . Buktikan bahwa $a\equiv 1(mod p^{n-1}).$
\end{example}
\begin{proof}
Jelas bahwa $p$ tidak membagi $a$.\\
Asumsikan $a\not\equiv 1(mod p^{n-1})$ , maka $a-1 \not\equiv 0(mod p^{n-1})$  .\\
Akibatnya , $v_p(a-1)<n-1$\\
Maka , dengan teorema LTE ,
\begin{equation} 
    \begin{split}
    v_p(a^p-1) & = v_p(a-1) +v_p(p)\\
               & < n-1 +1 \\
               & = n
    \end{split}
\end{equation}
Akibatnya $a^p-1 \not\equiv 0 \pmod{p^n}$ , kontradiksi.\\
Maka , $a \equiv 1 \pmod{p^{n-1}}$ .
\end{proof}

\section{Soal-Soal}
\begin{enumerate}
    \item Sebut suatu bilangan bulat sebagai \textit{square-free} jika bilangan tersebut tidak memiliki faktor bilangan kuadrat sempurna selain 1 . Jika $a>2$ adalah bilangan genap , buktikan bahwa bilangan $a^{a-1}-1$ bukan merupakan \textit{square-free}.
    \item Carilah semua solusi $(a,p,k)$ ,  dengan $a,k$ adalah bilangan bulat dan $p$ adalah bilangan prima , dari persamaan $$a^p-1=p^k$$
    \item Carilah semua pasangan bilangan bulat $(a,b)$ dengan $a,b>1$ sehingga $b^a|a^b-1$
    \item (Iran Second Round 2008) Tunjukkan bahwa satu-satunya nilai $a$ sehingga $4(a^n+1)$ merupakan bilangan kubik untuk setiap bilangan asli $n$ adalah 1
    \item Tunjukkan bahwa $a^n-b^n$ memiliki pembagi prima yang bukan merupakan pembagi dari $a-b$
    \item (Irlandia 1996) Misalkan $p$ merupakan bilangan prima dan $a,n$ merupakan bilangan asli . Buktikan bahwa jika $2^p+3^p=a^n$ maka $n=1$
    \item Carilah semua bilangan asli $n$  sehingga $$2^n|3^n-1$$
    \item (Rumania TST 1994) Misalkan $n$ adalah bilangan ganjil positif . Buktikan bahwa $((n-1)^n+1)^2$ membagi $n(n-1)^{(n-1)^n+1}+n$
    \item Misalkan $p\geq 5$ adalah bilangan prima . Carilah nilai maksimal dari bilangan asli $k$ sehingga $$p^k|(p-2)^{2(p-1)}-(p-4)^{(p-1)}$$
    \item (IMO SL 1991) Carilah nilai bilangan asli $k$ terbesar sehingga $1991^k$ habis membagi $$1990^{1991^{1992}}+1992^{1991^{1990}}$$
    \item Carilah semua bilangan asli $a$ sehingga $3^a$ membagi $5^a+1$
    \item (Bulgaria 1997) Buktikan bahwa jika $n$ adalah bilangan asli dan bilangan $3^n-2^n$ merupakan suatu bilangan prima berpangkat , maka $n$ adalah bilangan prima 
    \item Carilah semua solusi bilangan asli dari persamaan $$(n-1)!+1=n^m$$
    \item (Rusia 1996) Misalkan $x,y,n,p,k$ adalah bilangan asli dengan $n$ ganjil dan $p$ prima ganjil . Buktikan bahwa jika $x^n+y^n=p^k$ , maka $n=p^m$ untuk suatu bilangan asli $m$
    \item (Romanian Junior TST 2008) Misalkan $p$ adalah bilangan prima , $p \neq 3$ dan $a,b$ addalah bilangan bulat sehingga $p|a+b$ dan $p^2|a^3+b^3$ . Buktikan bahwa $p^2|a+b$ atau $p^3|a^3+b^3$
    \item (IMO 1990)Temukan semua bilangan asli $n$ sehingga $$\dfrac{2^n+1}{n^2}$$ adalah bilangan bulat .
\end{enumerate}
\end{document}